    %%%%%%%%%%%%%%%%%%%%%%%%%%%%%%%%%%%%%%%%%
% Medium Length Graduate Curriculum Vitae
% LaTeX Template
% Version 1.1 (9/12/12)
%
% This template has been downloaded from:
% http://www.LaTeXTemplates.com
%
% Original author:
% Rensselaer Polytechnic Institute (http://www.rpi.edu/dept/arc/training/latex/resumes/)
%
% Important note:
% This template requires the res.cls file to be in the same directory as the
% .tex file. The res.cls file provides the resume style used for structuring the
% document.
%
%%%%%%%%%%%%%%%%%%%%%%%%%%%%%%%%%%%%%%%%%

%----------------------------------------------------------------------------------------
%	PACKAGES AND OTHER DOCUMENT CONFIGURATIONS
%----------------------------------------------------------------------------------------

\documentclass[margin, 12pt]{res} % Use the res.cls style, the font size can be changed to 11pt or 12pt here
\usepackage{helvet} % Default font is the helvetica postscript font
%\usepackage{newcent} % To change the default font to the new century schoolbook postscript font uncomment this line and comment the one above
\usepackage{hyperref}
\usepackage[symbol]{footmisc}
\hypersetup{
	colorlinks = true,
	urlcolor = {blue}
}
\usepackage[inline]{enumitem}
\makeatletter
% This command ignores the optional argument for itemize and enumerate lists
\newcommand{\inlineitem}[1][]{%
	\ifnum\enit@type=\tw@
	{\descriptionlabel{#1}}
	\hspace{\labelsep}%
	\else
	\ifnum\enit@type=\z@
	\refstepcounter{\@listctr}\fi
	\quad\@itemlabel\hspace{\labelsep}%
	\fi}
\makeatother

\setlength{\textwidth}{5.1in} % Text width of the document

\begin{document}
	
	%----------------------------------------------------------------------------------------
	%	NAME AND ADDRESS SECTION
	%----------------------------------------------------------------------------------------
	
	\moveleft.5\hoffset\centerline{\large\bf Lorenz Gschwent} % Your name at the top
	
	\moveleft\hoffset\vbox{\hrule width\resumewidth height 1pt}\smallskip % Horizontal line after name; adjust line thickness by changing the '1pt'

	\moveleft.5\hoffset\centerline{lorenz.gschwent@uni-due.de}
	\moveleft.5\hoffset\centerline{\href{https://lorenzgschwent.github.io/}{Website}}
	\moveleft.5\hoffset\centerline{\href{https://github.com/lorenzgschwent/}{GitHub}}

	
	%----------------------------------------------------------------------------------------
	
	\begin{resume}
		
		%----------------------------------------------------------------------------------------
		%	RESERACH INTERESTS SECTION
		%----------------------------------------------------------------------------------------
		
		%\section{RESEARCH INTERESTS}
		%Regional Economics, Political Economy
		%\hspace{.1cm}
		%----------------------------------------------------------------------------------------
		%	EDUCATION SECTION
		%----------------------------------------------------------------------------------------
		
		\section{EDUCATION}
		
		{\sl PhD} in Economics \hfill since October 2022 \\
		University of Duisburg-Essen, \\
		RTG Regional Disparities and Economic Policy \\
		Supervisor: Tobias Seidel \vspace{0.2cm}\\
		{\sl Master of Science} in Economics \hfill 2019 - 2021 \\
		University of Vienna, Academic track \vspace{0.2cm}\\
		{\sl Erasmus Semester} \hfill September 2020 - January 2021 \\
		Aix-Marseille School of Economics, \\
		``Econometrics, Big Data, Statistics" track \vspace{0.2cm}\\
		{\sl Bachelor of Science} in Economics \hfill 2015 - 2019 \\
		University of Vienna \vspace{0.2cm}\\
		{\sl Bakkalaureat of Philosophy} in Communication Science \hfill 2014 - 2018 \\
		University of Vienna

		%----------------------------------------------------------------------------------------
		%	PROJECTS
		%----------------------------------------------------------------------------------------
		\section{PUBLIC- ATIONS}
		``The Rise of Health Economics: Transforming the Landscape of Economic Research". With Björn Hammarfelt, Martin Karlsson, and Mathias Kifmann. \textit{Health Economics}.
		
		\section{CURRENT PROJECTS}
		
		``Selection in the Patent Lottery" \\
		{\sl EPIP Young Scholar Award 2025}
		
		``Censorship in Democracy". With Marcel Caesmann and Matteo Grigoletto. \href{https://www.zora.uzh.ch/id/eprint/260103/}{Working Paper}. (Submitted). \\
		{\sl NOeG Young Economist Award 2023}, {\sl SSES Young Economist Award 2025}
		
		%----------------------------------------------------------------------------------------
		%	PRESENTATIONS SECTION	
		%----------------------------------------------------------------------------------------
		\section{PRESENT- ATIONS}
		{\sl2026}: DRUID Academy, ZEW PRICE Seminar, ICSSI Boulder, Melbourne Institute Brownbag %\footnote[1]{ scheduled}
		
		{\sl2025}: 18$^{th}$ RGS Doctoral Conference, 2$^{nd}$ Duke Strategy PhD Conference, BSE Summer Forum Economics of Science and Innovation, REGIS Summer School, EPIP Annual Conference, 8$^{th}$ RISE Workshop %\footnote[1]{ scheduled} %\footnotemark[1]
		
		{\sl2024}: 7$^{th}$ Conference on the Political Economy of Democracy and Dictatorship, 17$^{th}$ RGS Doctoral Conference
		
		{\sl2023}: CINCH Research Agenda for Ukraine Workshop (online), 5$^{th}$ Monash-Warwick-Zurich Text-as-Data Workshop (online), NOeG (Annual Conference of the Austrian Economic Association), Digital Democracy Workshop DigDemLab, 17$^{th}$ CESifo Workshop on Political Economy, 2$^{nd}$ CESifo/ifo Junior Workshop on Big Data
		%\footnote{ scheduled}\addtocounter{footnote}{-1}
		%\footnotemark\addtocounter{footnote}{-1}
		
		%----------------------------------------------------------------------------------------
		%	PUBLICATIONS SECTION
		%----------------------------------------------------------------------------------------
		%\section{PUBLI- CATIONS}
		\section{OTHER WORK}
		
		\href{https://wiiw.ac.at/digital-tasks-and-ict-capital-methodologies-and-data-p-6429.html}{Digital Tasks and ICT Capital: Methodologies and Data} (with Dario Guarascio, Stefan Jestl, Alireza Sabouniha, and Roman Stöllinger), wiiw Statistical Report No. 11, January 2023
		
		% Narrative-Character-Role package
		
		%----------------------------------------------------------------------------------------
		%	AWARDS SECTION	
		%----------------------------------------------------------------------------------------
		%\section{AWARDS}
		%Young Economist Award 2023
		
		%----------------------------------------------------------------------------------------
		%	TEACHING SECTION
		%----------------------------------------------------------------------------------------
		\section{TEACHING}
		{\sl Teaching Assistant} for Empirical Economics (yearly, since summer term 2023)\\
		{\sl Supervision:} Bachelor Seminar in Economics (since winter term 2022/'23), Master Seminar Applied Economic Research (summer term 2023)
		
		%----------------------------------------------------------------------------------------
		%	VISITING
		%----------------------------------------------------------------------------------------
		\section{VISITING}
		{\sl ZEW Mannheim} \hfill February/March 2026 \\
		{\sl Melbourne Institute} \hfill September-November 2026
		
		%----------------------------------------------------------------------------------------
		%	PREVIOUS POSITIONS SECTION
		%----------------------------------------------------------------------------------------
		\section{PREVIOUS POSITIONS}
		{\sl Research Associate} \hfill October 2021 - August 2022\\ Wyss Academy for Nature at the University of Bern\\
		Pre-Doctoral Research Assistant with Kai Gehring \vspace{0.2cm} \\
		{\sl Part-time Research Assistant (50\%)} \hfill July 2021 - September 2021\\
		The Vienna Institute for International Economic Studies (wiiw)
		%Research assistance on the project \href{https://wiiw.ac.at/comparative-advantage-in-the-digital-era-new-insights-into-trade-in-digital-tasks-and-ict-capital-pj-244.html}{``Comparative Advantage in the Digital Era: New insights into trade in digital tasks and ICT capital"} 
		
		%----------------------------------------------------------------------------------------
		%	COMPUTER SKILLS SECTION
		%----------------------------------------------------------------------------------------
		\section{COMPUTER \\ SKILLS} 
		
		Python, STATA, R, Matlab, SQL, Bash Script, LaTeX, and Microsoft Office
		
		
		%----------------------------------------------------------------------------------------
		%	LANGUAGES
		%----------------------------------------------------------------------------------------
		
		\section{LANGUAGES}
		
		German (native), English

		
	\end{resume}
\end{document}